\section{总结}

\begin{frame}{主要结论}
  \begin{enumerate}
    \item 心理学版图灵测试将 LLM 评估从“是否正确”推进到“是否具有类人社会理性模式”。
    \item 系统支持通过结构化配置扩展新 LLM、新版本与新问卷，并自动生成问卷特定指标。
    \item 初步结果显示，当前中国 LLM 与人类存在系统性差异，尤其是在决策前提任务的社会线索敏感性上。
    \item LLM 之间存在显著不一致；思考模式不是简单优势，而是会改变行为误差结构。
  \end{enumerate}
\end{frame}

\begin{frame}{对 LLM 评估的启示}
  \begin{block}{方法层面}
    传统准确率基准很难捕捉亲缘关系、规模、收益变化、身份线索和自我相关前提对判断的影响。
    社会理性基准可以把这些心理学结构转化为可复现的模型画像。
  \end{block}

  \vspace{0.6em}
  \begin{block}{模型层面}
    更强的推理能力或更高版本号不必然带来更高的类人一致性。
    对齐需要同时看具体行为结果与跨条件的心理价值结构。
  \end{block}

  \vspace{0.6em}
  \centering
  \Large\metric{类人判断能力需要被“心理学地”测量。}
\end{frame}
